% 研究方法：實驗一包含 Baseline 與集成；實驗二為 FS 集成。
\chapter{研究方法}\label{chap:method}
\begingroup
\setlength{\emergencystretch}{2em}
\clubpenalty=10000
\widowpenalty=10000
\displaywidowpenalty=10000

\section{研究架構}\label{sec:method-framework}

本研究以企業破產預測為研究主題，探討訓練資料期間、分類器、採樣策略、集成方法與特徵選取對分類效能之影響。研究考量公司年度財務資料的類別不平衡特性，以及不同年度資料可能存在的分布差異，透過歷史資料與近期資料所建立之模型，分析跨期間資訊在破產類別辨識上的應用方式。

本研究分為實驗一與實驗二，整體架構如圖~\ref{fig:method-overview} 所示。實驗一包含 Baseline 分類器比較與 Old／New 集成比較：先比較分類器、訓練資料期間與採樣策略，再固定分類器種類建立六模型池，比較靜態集成、動態分類器選擇與動態集成選擇。實驗二延續實驗一的集成設定，加入不同特徵選取方法，重新建立模型池並評估 FS 對集成分類效能的影響。

表~\ref{tab:method-stages} 整理兩項實驗的比較重點。實驗一的 Baseline 比較與集成比較為同一實驗內的兩個部分；實驗二則以未進行 FS 的集成結果為對照。各部分依對應的資料分區與特徵設定重新訓練模型。

\begingroup
\small
\setstretch{1.15}
\setlength{\tabcolsep}{4pt}
\renewcommand{\arraystretch}{1.25}
\setlength{\LTleft}{0pt plus 1fill}
\setlength{\LTright}{0pt plus 1fill}
\begin{longtable}{@{}>{\raggedright\arraybackslash}p{\dimexpr 0.18\linewidth-1.0800\tabcolsep\relax}>{\raggedright\arraybackslash}p{\dimexpr 0.28\linewidth-1.6800\tabcolsep\relax}>{\raggedright\arraybackslash}p{\dimexpr 0.28\linewidth-1.6800\tabcolsep\relax}>{\raggedright\arraybackslash}p{\dimexpr 0.26\linewidth-1.5600\tabcolsep\relax}@{}}
\caption{兩項實驗之研究架構與比較目的}\label{tab:method-stages}\\
\toprule
\textbf{實驗} & \textbf{主要變動因素} & \textbf{比較目的} & \textbf{實驗間之銜接} \\
\midrule
\endfirsthead
\multicolumn{4}{c}{表~\thetable（續）}\\
\toprule
\textbf{實驗} & \textbf{主要變動因素} & \textbf{比較目的} & \textbf{實驗間之銜接} \\
\midrule
\endhead
\midrule
\multicolumn{4}{r}{續下頁}\\
\endfoot
\bottomrule
\endlastfoot
實驗一：Baseline 與集成比較 & 分類器、訓練期間、採樣方式及模型組成 & 建立單模型基準，分析 Old／New 模型的集成方式 & 提供實驗二未進行 FS 的集成對照 \\
\addlinespace[2pt]
實驗二：FS 集成比較 & 特徵選取方法及保留數 & 比較加入 FS 後的集成分類效能 & 沿用實驗一集成設定，進行同條件配對比較 \\*
\end{longtable}
\endgroup
\begin{figure}[htbp]
\centering
\begingroup
\small
\setstretch{1.12}
\setlength{\fboxsep}{8pt}
\begin{tabular}{@{}c@{}}
\fbox{\parbox{0.88\linewidth}{\centering 公司年度財務資料\\
定義破產正類；依年度劃分擬合、驗證與測試用途}}\\[3pt]
\(\Downarrow\)\\[3pt]
\fbox{\parbox{0.88\linewidth}{\centering\textbf{實驗一：Baseline 與集成分類器比較}\\[4pt]
Baseline：分類器、Old／New／Retrain、採樣策略之比較\\
\(\Downarrow\)\\
固定分類器種類，建立 Old 三模型與 New 三模型\\
比較靜態集成、DCS 與 DES}}\\[3pt]
\(\Downarrow\)\\[3pt]
\fbox{\parbox{0.88\linewidth}{\centering\textbf{實驗二：特徵選取集成分類器比較}\\[4pt]
沿用集成比較之資料分區、分類器種類與採樣方式\\
以 Old 擬合資料選取特徵，再採樣並重新建立六模型池\\
比較加入 FS 前後的固定方法與驗證選模程序}}\\[3pt]
\(\Downarrow\)\\[3pt]
\fbox{\parbox{0.88\linewidth}{\centering 驗證決策固定後，評估 2015--2018 年測試資料\\
主要報告 AUC、F1-score 與 Recall}}
\end{tabular}
\endgroup
\caption{本研究兩項實驗之整體架構}\label{fig:method-overview}
\end{figure}

本研究以公司年度財務紀錄為觀測單位，使用資料集提供的公司最終存續或破產狀態作為分類標籤，並將破產類別設為正類。模型輸入為財務特徵，公司識別碼、年度及狀態標籤不納入預測特徵。由於同一公司可有多年度紀錄，年份切分後的訓練與測試資料仍可能包含相同公司，因此，本研究評估的是較晚年度公司財務紀錄的分類表現，並不等同於完全未見公司的泛化能力。原始標籤亦未建立明確的事件預測期間，故結果以回溯性公司狀態分類解釋，不據此推論模型已具備提前一年預測破產事件的能力。

最終測試以 ROC 曲線下面積（Area Under the ROC Curve, AUC）為主要報告指標，並以 F1-score 與 Recall 補充破產類別的辨識表現。實驗一之集成比較與實驗二另以平均精確率（Average Precision, AP）選擇靜態集成候選，並依驗證資料的 F1-score 決定分類門檻。候選選擇、門檻決定與測試報告分別具有不同用途，其規則於後續各節說明。

\clearpage
\section{實驗一：Baseline 與集成分類器比較}\label{sec:method-experiment1}

實驗一包含 Baseline 分類器比較與 Old／New 集成比較。前者建立單一分類器基準並分析訓練資料期間與採樣策略；後者固定分類器種類，分析六模型池的靜態與動態整合方式。兩部分的資料分區分別說明，避免將其分數差直接解讀為集成帶來的改善。

\subsection{Baseline 分類器比較與選擇}\label{sec:method-phase1}

Baseline 部分比較不同分類器在訓練資料期間與採樣設定改變時的分類表現，建立實驗一的單模型基準。分類器、資料期間、訓練策略與採樣方式構成平行比較設定，其流程如圖~\ref{fig:method-baseline} 所示。
\begin{figure}[htbp]
\centering
\includegraphics[width=\linewidth]{figures/experiment1-baseline.png}
\caption{實驗一之 Baseline 分類器比較架構}\label{fig:method-baseline}
\end{figure}



\subsubsection{實驗資料與年份切分}

本研究使用 1999 年至 2018 年之公司年度財務資料，將 1999 年至 2014 年設定為模型建立期間，2015 年至 2018 年固定作為測試期間。為分析歷史與近期資料範圍對分類結果的影響，本部分逐年移動 Old／New 分界，形成 15 組年份設定。

令 \(b\) 為 Old 期間的最後一年，則各組資料期間定義為：

\begin{equation}\label{eq:method-years}
\begin{aligned}
\mathrm{Old}(b)&=\{1999,\ldots,b\},\\
\mathrm{New}(b)&=\{b+1,\ldots,2014\},\\
b&\in\{1999,\ldots,2013\}.
\end{aligned}
\end{equation}

因此，最早一組為 Old 使用 1999 年、New 使用 2000 年至 2014 年；最晚一組為 Old 使用 1999 年至 2013 年、New 使用 2014 年。各組皆保留至少一年的 Old 與 New 資料。

在內層驗證方面，既有 Baseline 比較流程從對應訓練期間的各年份資料中抽取約 20\% 樣本作為驗證資料，其餘作為實際擬合資料。Old 與 New 分別依其期間建立擬合與驗證分區；Retrain 則合併兩側的擬合資料及驗證資料。前處理與採樣於切分後執行，驗證資料用於相關設定與門檻的決定。

此年內抽樣方式使擬合與驗證資料可能涵蓋相同年份，與實驗一之集成比較與實驗二完整保留 2014 年作為驗證期間的方式不同。因此，15 組設定主要用於觀察訓練期間敏感度；跨比較設定分數的比較另須考量驗證制度與資料可用範圍的差異。

\subsubsection{分類器與訓練策略}

候選分類器涵蓋傳統機器學習與以神經網路為基礎的表格模型，包括 Logistic Regression、Random Forest、XGBoost、LightGBM、TabM 與 TabICL。主要超參數設定列於本章，補充訓練控制與後端版本說明見附錄~\ref{app:experimental-details}。

本部分比較三種訓練策略。Old 僅使用歷史期間的擬合資料，評估歷史資訊對較晚年度樣本的適用程度；New 使用近期期間的擬合資料，評估近期資訊所建立的分類能力；Retrain 合併 Old 與 New 的擬合資料，從頭訓練單一模型，作為使用完整可用訓練期間的比較基準。

本章以 New 表示使用近期資料建立的模型分支。模型是否延續舊模型權重屬於不同的訓練程序，相關名稱與版本的區分見附錄~\ref{app:experimental-details}。

\subsubsection{資料前處理與採樣策略}

資料前處理包含依各分類器流程執行的遺失值處理與數值縮放。需由資料估計的參數僅使用對應的擬合資料計算，再將固定轉換套用於驗證與測試資料。不同前處理設定之比較界線見附錄~\ref{app:experimental-details}。

採樣設定包含不採樣（None）、低採樣（Under-sampling）、過採樣（Over-sampling）與混合採樣（Hybrid Sampling）。不採樣保留原始類別分布；其餘設定分別透過樣本刪減、少數類別樣本生成，或兩者結合，改變模型實際使用的訓練資料。採樣僅作用於擬合資料，驗證與測試資料均保留原始類別分布。

本研究採用的代表方法為 TomekLinks、ADASYN 與 SMOTEENN。TomekLinks 依近鄰關係移除類別邊界附近的指定樣本，其作用偏向邊界清理，不能視為固定將兩類樣本調整至相同比例。各方法的參數列於實驗設定，補充版本說明見附錄。

\subsubsection{分類器比較與選擇方式}

本部分的比較單位為「分類器 \ensuremath{\times} 年份邊界 \ensuremath{\times} 訓練策略 \ensuremath{\times} 採樣方式」。研究設計要求分類器種類、超參數與分類門檻依訓練及驗證資料決定，固定後再於 2015 年至 2018 年資料計算 AUC、F1-score 與 Recall。

分類器種類的選擇與單一設定內的超參數調整分別記錄。XGBoost、LightGBM、Random Forest、Logistic Regression 與 TabM 的調參流程以驗證 AUC 選擇參數，候選集合與訓練控制列於第~\ref{sec:method-phase1-params}~節；後續分類器種類的選擇，仍須另行交代如何彙整不同年份邊界、訓練策略與採樣設定的驗證結果。此選擇程序不得依測試結果反覆改定；若僅為既有測試結果的描述性比較，則不將其表述為已完成的獨立驗證選模。

由於 15 組設定共用測試期間，且訓練資料高度重疊，各組分數具有相依性。分析時以期間敏感度與描述性差異為主，不將其視為 15 次彼此獨立的驗證。分類器選擇與固定設定的比較限制見附錄~\ref{app:experimental-details}。

\clearpage
\subsection{Old／New 集成分類器比較}\label{sec:method-phase2}

實驗一的集成部分固定基礎分類器種類，分析不同資料期間與採樣方式所形成模型的整合表現。本部分使用 XGBoost 建立 Old／New 模型池，比較靜態集成、DCS 與 DES，架構如圖~\ref{fig:method-ensemble} 所示。
\begin{figure}[htbp]
\centering
\includegraphics[width=\linewidth]{figures/experiment1-ensemble.png}
\caption{實驗一之 Old／New 集成分類器架構}\label{fig:method-ensemble}
\par\smallskip{\footnotesize 圖中 Validation 為驗證用途的概括表示；實際分為互不重疊的 DSEL 與 Calibration。前者估計局部分類能力，後者選擇靜態候選及分類門檻。\par}
\end{figure}



\subsubsection{實驗資料切分}

本部分以 1999 年至 2011 年為 Old 訓練期間，2012 年至 2014 年為名義上的 New 營運期間，並固定 2015 年至 2018 年為測試期間。其中，2014 年完整保留供驗證使用，因此 New 模型實際以 2012 年至 2013 年資料進行擬合。

2014 年資料以固定隨機種子 42，依類別標籤分層切分為等量且互不重疊的動態選擇資料集（Dynamic Selection Dataset, DSEL）與校準資料集（Calibration）。DSEL 提供動態方法估計局部分類能力；Calibration 用於靜態候選選擇與分類門檻決定，兩者均不加入基礎分類器的擬合資料。

\begingroup
\footnotesize
\setstretch{1.15}
\setlength{\tabcolsep}{4pt}
\renewcommand{\arraystretch}{1.15}
\setlength{\LTleft}{0pt plus 1fill}
\setlength{\LTright}{0pt plus 1fill}
\begin{longtable}{@{}>{\raggedright\arraybackslash}p{\dimexpr 0.15\linewidth-1.2000\tabcolsep\relax}>{\raggedright\arraybackslash}p{\dimexpr 0.16\linewidth-1.2800\tabcolsep\relax}>{\raggedright\arraybackslash}p{\dimexpr 0.17\linewidth-1.3600\tabcolsep\relax}>{\raggedright\arraybackslash}p{\dimexpr 0.15\linewidth-1.2000\tabcolsep\relax}>{\raggedright\arraybackslash}p{\dimexpr 0.37\linewidth-2.9600\tabcolsep\relax}@{}}
\caption{集成比較之固定年份資料分區}\label{tab:method-partitions}\\
\toprule
\textbf{分區} & \textbf{年份} & \textbf{公司年度紀錄數} & \textbf{正類紀錄數} & \textbf{主要用途} \\
\midrule
\endfirsthead
\multicolumn{5}{c}{表~\thetable（續）}\\
\toprule
\textbf{分區} & \textbf{年份} & \textbf{公司年度紀錄數} & \textbf{正類紀錄數} & \textbf{主要用途} \\
\midrule
\endhead
\midrule
\multicolumn{5}{r}{續下頁}\\
\endfoot
\bottomrule
\endlastfoot
Old-fit & 1999–2011 & 55,927 & 4,442 & Old 模型訓練；前處理與特徵選取擬合 \\
\addlinespace[2pt]
New-fit & 2012–2013 & 6,969 & 349 & New 模型訓練 \\
\addlinespace[2pt]
DSEL & 2014 & 1,752 & 71 & 動態選擇的局部分類能力估計 \\
\addlinespace[2pt]
Calibration & 2014 & 1,752 & 71 & 靜態候選選擇與分類門檻決定 \\
\addlinespace[2pt]
Test & 2015–2018 & 12,282 & 287 & 設定固定後的最終評估 \\
\addlinespace[2pt]
合計 & 1999–2018 & 78,682 & 5,220 & 五個分區在資料列上互斥 \\*
\end{longtable}
\endgroup

資料來源：本研究實驗資料分區彙整。表中數值為公司年度紀錄數，並非各分區獨立公司數。

本研究沿用 Calibration 作為分區名稱，但其程序為候選與門檻選擇，未另執行機率校準。DSEL 與 Calibration 均取自 2014 年，資料未提供更細的年內可得時間，因此不將兩者描述為具有年內先後順序的驗證批次。

\subsubsection{基礎模型池建立}

本部分先於 Old-fit 估計各財務特徵的平均數，以進行遺失值補值，再以補值後的 Old-fit 擬合標準化參數。相同轉換一致套用於 New-fit、DSEL、Calibration 與 Test，使六個基礎模型及動態方法使用共同的特徵表示。

完成前處理後，Old-fit 與 New-fit 各自執行 TomekLinks、ADASYN 及 SMOTEENN，並分別重新訓練分類器，形成表~\ref{tab:method-pool} 所示的六模型池。實驗一之集成比較使用完整 18 個財務特徵；分類器本身不另加入類別權重，以採樣方式作為本部分的類別不平衡處理因素。

\begingroup
\small
\setstretch{1.15}
\setlength{\tabcolsep}{4pt}
\renewcommand{\arraystretch}{1.25}
\setlength{\LTleft}{0pt plus 1fill}
\setlength{\LTright}{0pt plus 1fill}
\begin{longtable}{@{}>{\raggedright\arraybackslash}p{\dimexpr 0.28\linewidth-1.1200\tabcolsep\relax}>{\raggedright\arraybackslash}p{\dimexpr 0.25\linewidth-1.0000\tabcolsep\relax}>{\raggedright\arraybackslash}p{\dimexpr 0.47\linewidth-1.8800\tabcolsep\relax}@{}}
\caption{Old／New 六模型池組成}\label{tab:method-pool}\\
\toprule
\textbf{模型名稱} & \textbf{擬合資料} & \textbf{採樣方法} \\
\midrule
\endfirsthead
\multicolumn{3}{c}{表~\thetable（續）}\\
\toprule
\textbf{模型名稱} & \textbf{擬合資料} & \textbf{採樣方法} \\
\midrule
\endhead
\midrule
\multicolumn{3}{r}{續下頁}\\
\endfoot
\bottomrule
\endlastfoot
Old\_under & Old-fit & TomekLinks \\
\addlinespace[2pt]
Old\_over & Old-fit & ADASYN \\
\addlinespace[2pt]
Old\_hybrid & Old-fit & SMOTEENN \\
\addlinespace[2pt]
New\_under & New-fit & TomekLinks \\
\addlinespace[2pt]
New\_over & New-fit & ADASYN \\
\addlinespace[2pt]
New\_hybrid & New-fit & SMOTEENN \\*
\end{longtable}
\endgroup

六個分類器使用相同的模型種類，透過訓練期間與採樣方式產生不同模型。此設計提供模型差異的來源，但不預設其預測一定具有互補性；互補資訊能否改善分類效能，仍須由驗證與測試結果觀察。

\subsubsection{靜態集成方法}

靜態集成於推論前固定成員組合，並以成員對破產正類的預測機率進行等權平均。令模型池中的三個 Old 模型集合為 \(\mathcal{O}\)，三個 New 模型集合為 \(\mathcal{N}\)，則混合候選子集須滿足：

\begin{equation}\label{eq:method-subsets}
S\subseteq\mathcal{O}\cup\mathcal{N},\qquad
S\cap\mathcal{O}\ne\varnothing,\qquad
S\cap\mathcal{N}\ne\varnothing.
\end{equation}

Old 與 New 各有 \(2^3-1\) 個非空子集，因此混合組合數為：

\begin{equation}\label{eq:method-count}
(2^3-1)(2^3-1)=49.
\end{equation}

\begingroup
\small
\setstretch{1.15}
\setlength{\tabcolsep}{4pt}
\renewcommand{\arraystretch}{1.25}
\setlength{\LTleft}{0pt plus 1fill}
\setlength{\LTright}{0pt plus 1fill}
\begin{longtable}{@{}>{\raggedright\arraybackslash}p{\dimexpr 0.55\linewidth-1.1000\tabcolsep\relax}>{\raggedright\arraybackslash}p{\dimexpr 0.45\linewidth-0.9000\tabcolsep\relax}@{}}
\caption{不同模型數量的 Old／New 混合靜態組合}\label{tab:method-combinations}\\
\toprule
\textbf{集成模型數量} & \textbf{混合組合數} \\
\midrule
\endfirsthead
\multicolumn{2}{c}{表~\thetable（續）}\\
\toprule
\textbf{集成模型數量} & \textbf{混合組合數} \\
\midrule
\endhead
\midrule
\multicolumn{2}{r}{續下頁}\\
\endfoot
\bottomrule
\endlastfoot
2 & 9 \\
\addlinespace[2pt]
3 & 18 \\
\addlinespace[2pt]
4 & 15 \\
\addlinespace[2pt]
5 & 6 \\
\addlinespace[2pt]
6 & 1 \\
\addlinespace[2pt]
合計 & 49 \\*
\end{longtable}
\endgroup

對任一候選子集 \(S\)，集成機率定義為：

\begin{equation}\label{eq:method-probability}
\hat p_S(x)=\frac{1}{|S|}\sum_{i\in S}\hat p_i(x),
\end{equation}

其中，\(\hat p_i(x)\) 為第 \(i\) 個分類器對樣本 \(x\) 預測為破產的機率，\(|S|\) 為該子集的模型數。

除混合候選外，另保留六個單模型，以及三個 Old 模型的等權集成 Old3、三個 New 模型的等權集成 New3，作為比較基準。完整六模型的等權集成稱為 Equal6，其成員組合已包含在 49 種混合候選之中，不額外增加混合組合數。Old3 與 New3 為單一期間控制組，均不納入 49 種跨期間混合候選的選擇範圍。

\subsubsection{動態分類器與集成選擇}

動態選擇依待預測樣本的鄰近區域評估基礎模型的分類能力。本研究以完整六模型池為候選來源，使用獨立 DSEL 建立局部能力估計，並比較動態分類器選擇（Dynamic Classifier Selection, DCS）與動態集成選擇（Dynamic Ensemble Selection, DES）。

DCS 比較 Overall Local Accuracy（OLA）與 Local Class Accuracy（LCA）。OLA 以分類器在樣本鄰近區域的整體正確率估計其能力；LCA 則依各分類器對待預測樣本所判定的類別，估計相應的局部分類能力。兩者原則上選擇局部能力最高的單一分類器，平手時依固定模型順序選擇。

DES 比較 k-Nearest Oracles Eliminate（KNORA-E）與 k-Nearest Oracles Union（KNORA-U）。KNORA-E 尋找能正確分類鄰近區域所有樣本的分類器，必要時縮小鄰域；KNORA-U 則選取至少能正確分類一個鄰近樣本的分類器，並依其局部正確分類次數形成貢獻權重。兩者依所選分類器產生最終集成輸出。

本研究使用 DESlib 0.3.7 的標準實作，鄰居數固定為 \(k=7\)；DES 採機率軟投票，DCS 使用 {\normalfont best} 選擇規則，並關閉額外的動態剪枝與樣本難度分流。當所有基礎分類器對某樣本的類別判斷一致時，沿用套件直接平均機率的處理。因此，DCS 的單分類器選擇描述適用於需要進行局部選擇的情況，不能推定每個樣本的輸出皆來自單一模型。

上述方法的標準定義分別參見 DESlib 的 OLA\autocite{deslib-ola-docs}、LCA\autocite{deslib-lca-docs}、KNORA-E\autocite{deslib-knorae-docs} 與 KNORA-U\autocite{deslib-knorau-docs} 文件；本研究的實際版本與開關設定以正式執行紀錄為準。

\subsubsection{集成候選選擇與分類門檻}

靜態候選與分類門檻均由 Calibration 決定。首先，對 49 種混合靜態候選計算 Calibration AP，選出 AP 最高者，記為 ValidationSelectedStatic。若 AP 相同，依固定的方法名稱字典序決定，以維持選擇結果的可重現性。保留所有候選的測試表現可供組成分析，但不再依 Test 排名回選候選。

其次，各單模型、控制組、靜態候選與動態方法，均在 Calibration 上搜尋 F1-score 最高的分類門檻。門檻候選集合為：

\begin{equation}\label{eq:method-threshold}
\mathcal{T}=\{0.01,0.02,\ldots,0.99\},\qquad
\tau_m^*=\underset{\tau\in\mathcal{T}}{\arg\max}\;
F1_{\mathrm{Calibration}}(m,\tau).
\end{equation}

若多個門檻得到相同的最高 F1-score，採用其中最小者。分類時以 \(\hat p_m(x)\ge\tau_m^*\) 判定為正類；ValidationSelectedStatic 沿用其所選候選的門檻。

候選與門檻固定後，於 Test 計算整體及各年度的 AUC、F1-score 與 Recall。各方法在四個測試年度均沿用相同的 Calibration 門檻，不利用年度測試標籤重新調整。

Retrain 在研究架構中定位為 Old-fit 與 New-fit 合併訓練的獨立基準，不納入六模型池。其比較須另以相同資料可用範圍、採樣及驗證規則建立；其完成範圍與同協定比較限制詳見附錄~\ref{app:experimental-details}。

\clearpage
\section{實驗二：特徵選取集成分類器比較}\label{sec:method-phase3}

實驗二延續實驗一的集成比較設定，固定資料分區、分類器種類、採樣方式與集成方法，在模型訓練前加入特徵選取，架構如圖~\ref{fig:method-fs} 所示。本實驗比較不同特徵集合下的集成分類效能，不預設特徵刪減必然改善結果。
\begin{figure}[htbp]
\centering
\includegraphics[width=\linewidth]{figures/experiment2-fs-ensemble.png}
\caption{實驗二之特徵選取集成分類器架構}\label{fig:method-fs}
\par\smallskip{\footnotesize 各 FS 方法平行執行；SHAP 為模型解釋後的特徵排序，其定位另於正文說明。DSEL 與 Calibration 維持實驗一集成部分的分配。\par}
\end{figure}



\subsection{特徵選取方法}

特徵選取可依機制概分為過濾式（Filter）、包裝式（Wrapper）及嵌入式（Embedded）。本研究的具體設定為互資訊（Mutual Information, MI）、分類與迴歸樹（Classification and Regression Tree, CART）、遞迴特徵消除（Recursive Feature Elimination, RFE），以及以 SHAP 貢獻值排序的特徵選取。各方法獨立執行，不依序串接。

\begingroup
\small
\setstretch{1.15}
\setlength{\tabcolsep}{4pt}
\renewcommand{\arraystretch}{1.25}
\setlength{\LTleft}{0pt plus 1fill}
\setlength{\LTright}{0pt plus 1fill}
\begin{longtable}{@{}>{\raggedright\arraybackslash}p{\dimexpr 0.15\linewidth-0.6000\tabcolsep\relax}>{\raggedright\arraybackslash}p{\dimexpr 0.28\linewidth-1.1200\tabcolsep\relax}>{\raggedright\arraybackslash}p{\dimexpr 0.57\linewidth-2.2800\tabcolsep\relax}@{}}
\caption{特徵選取方法與操作方式}\label{tab:method-fs}\\
\toprule
\textbf{方法} & \textbf{方法定位} & \textbf{特徵選取依據} \\
\midrule
\endfirsthead
\multicolumn{3}{c}{表~\thetable（續）}\\
\toprule
\textbf{方法} & \textbf{方法定位} & \textbf{特徵選取依據} \\
\midrule
\endhead
\midrule
\multicolumn{3}{r}{續下頁}\\
\endfoot
\bottomrule
\endlastfoot
MI & Filter & 計算各特徵與分類標籤的互資訊，依分數選取前 \(k_{\mathrm{FS}}\) 個特徵 \\
\addlinespace[2pt]
CART & 模型重要性篩選，對應 Embedded 類 & 以決策樹學得的特徵重要性排序，選取前 \(k_{\mathrm{FS}}\) 個特徵 \\
\addlinespace[2pt]
RFE & Wrapper & 以決策樹反覆擬合並移除特徵，直至保留指定數量 \\
\addlinespace[2pt]
SHAP 排序 & 模型解釋後的特徵篩選 & 以 LightGBM 計算訓練樣本的樹模型貢獻值，依平均絕對值排序選取特徵 \\*
\end{longtable}
\endgroup

MI 與 RFE 的機制可參照 scikit-learn 特徵選取文件\autocite{sklearn-fs-docs}。SHAP 設定使用 LightGBM 的樹模型貢獻值輸出，排除基準值欄位後計算各特徵的重要程度，其輸出定義見 LightGBM Booster 文件\autocite{lightgbm-booster-docs}。本研究另行交代 SHAP 的操作機制，不將其直接等同於模型訓練過程內完成的 Embedded 選取。

各方法分別設定保留比例 \(r=0.5\) 與 \(r=0.8\)，特徵數依下式決定：

\begin{equation}\label{eq:method-fs-count}
k_{\mathrm{FS}}=\max\{1,\lfloor d\times r\rfloor\},
\end{equation}

其中，\(d=18\) 為原始財務特徵數。因此，兩種設定分別保留 9 與 14 個特徵；80\% 為名義比例，實際保留數依向下取整規則計算。四種方法各有兩種保留設定，共形成八種特徵選取條件，另以完整 18 特徵作為對照。

\subsection{特徵選取與模型訓練流程}

本部分的處理順序固定為「資料分區 \ensuremath{\rightarrow} 前處理 \ensuremath{\rightarrow} 特徵選取 \ensuremath{\rightarrow} 採樣 \ensuremath{\rightarrow} 模型訓練」。補值與標準化參數先以 Old-fit 擬合，再於轉換後的 Old-fit 執行各特徵選取方法。所選特徵及欄位順序固定後，一致套用於 New-fit、DSEL、Calibration 與 Test。

特徵選取於採樣前執行，使其輸入資料保留 Old-fit 原始樣本組成。此安排不代表選取器不處理類別不平衡：CART、RFE 使用的決策樹及 SHAP 使用的 LightGBM 均設定平衡類別權重；MI 則依其互資訊估計程序執行。特徵選取器的設定與後續基礎分類器的設定分別記錄。

完成特徵選取後，Old-fit 與 New-fit 分別執行三種採樣，再訓練三個 Old 與三個 New 分類器。各特徵條件均重新建立六模型池，以確保模型於對應的輸入空間中學習。

同一特徵條件下，六模型、DSEL 與其他評估分區使用共同特徵集合；不同特徵條件則各自重新建立動態選擇器。由於特徵選取也會改變 DSEL 的距離計算空間，其影響可能同時來自基礎分類器與局部鄰域的改變。

\subsection{特徵選取後的集成比較}

各特徵條件均沿用實驗一之集成比較的 49 種混合靜態子集、Old3、New3、Equal6 與四種動態方法。DSEL 的資料列分配維持一致；Calibration 分別用於該條件的靜態候選選擇與門檻決定，Test 則使用相同樣本計算最終指標。

比較分為兩個層次。第一個層次固定方法與模型成員，於相同種子及測試樣本下比較特徵選取前後的 AUC、F1-score 與 Recall。第二個層次比較各條件的 ValidationSelectedStatic，評估加入特徵選取後，整套驗證選模程序的表現。後者允許 Calibration 選出不同的靜態成員組合，不能將其差異全部歸因於特徵刪減。

對任一固定方法 \(m\) 與指標 \(M\)，配對差值定義為：

\begin{equation}\label{eq:method-delta}
\Delta M_{s,m,f}
=M_{s,m,f}-M_{s,m,\mathrm{no\_fs}},
\end{equation}

其中，\(s\) 表示相同隨機種子，\(f\) 表示特徵選取條件。正值表示該指標高於未進行特徵選取的對照。即使模型成員固定，各條件的分類門檻仍依各自的 Calibration 決定，因此 F1-score 與 Recall 差異反映特徵選取、重新訓練及門檻選擇後的整體結果。除效能外，亦記錄所選特徵名稱、順序與數量，以保留結果追溯依據。

\section{實驗設定與參數}\label{sec:method-settings}

\subsection{資料與驗證設定}

本研究使用本地美國企業破產資料，共 78,682 筆公司年度紀錄，涵蓋 8,971 家公司與 18 個財務特徵；正類紀錄共 5,220 筆。同一公司可有多年度資料，因此紀錄數與公司數分別呈現。各部分的驗證方式整理於表~\ref{tab:method-split-summary}。

\begingroup
\small
\setstretch{1.15}
\setlength{\tabcolsep}{4pt}
\renewcommand{\arraystretch}{1.25}
\setlength{\LTleft}{0pt plus 1fill}
\setlength{\LTright}{0pt plus 1fill}
\begin{longtable}{@{}>{\raggedright\arraybackslash}p{\dimexpr 0.13\linewidth-0.7800\tabcolsep\relax}>{\raggedright\arraybackslash}p{\dimexpr 0.34\linewidth-2.0400\tabcolsep\relax}>{\raggedright\arraybackslash}p{\dimexpr 0.31\linewidth-1.8600\tabcolsep\relax}>{\raggedright\arraybackslash}p{\dimexpr 0.22\linewidth-1.3200\tabcolsep\relax}@{}}
\caption{兩項實驗資料與驗證設定摘要}\label{tab:method-split-summary}\\
\toprule
\textbf{比較部分} & \textbf{實際擬合資料} & \textbf{驗證用途} & \textbf{測試資料} \\
\midrule
\endfirsthead
\multicolumn{4}{c}{表~\thetable（續）}\\
\toprule
\textbf{比較部分} & \textbf{實際擬合資料} & \textbf{驗證用途} & \textbf{測試資料} \\
\midrule
\endhead
\midrule
\multicolumn{4}{r}{續下頁}\\
\endfoot
\bottomrule
\endlastfoot
實驗一：Baseline & 1999–2014 年內依 15 組邊界與策略設定，扣除年內驗證樣本 & 各訓練年份抽取約 20\%，供設定比較與門檻決定 & 2015–2018 年 \\
\addlinespace[2pt]
實驗一：集成 & Old-fit：1999–2011；New-fit：2012–2013 & 2014 年等分為 DSEL 與 Calibration & 2015–2018 年 \\
\addlinespace[2pt]
實驗二：FS 集成 & 同實驗一之集成比較；前處理及特徵選取僅擬合 Old-fit & 同實驗一之集成比較，固定資料列分配 & 2015–2018 年 \\*
\end{longtable}
\endgroup

\subsection{Baseline 分類器之超參數設定}\label{sec:method-phase1-params}

Baseline 分類器比較同時包含預設參數與超參數調整設定。XGBoost、LightGBM、Random Forest、Logistic Regression 與 TabM 先以擬合資料訓練候選，再依未採樣驗證資料的 ROC-AUC 選參，測試資料不參與決策。相關實驗紀錄、版本與比較界線見附錄~\ref{app:experimental-details}。

表~\ref{tab:method-ml-search} 與表~\ref{tab:method-tabm-search} 採「分類器、參數項目、候選值」列出本研究目前採用的候選設定。候選集合與單次搜尋預算分別定義；從集合中抽取有限組合，不代表窮舉全部排列。四種傳統分類器的搜尋預算預設為 48 組，可由執行參數覆寫。Logistic Regression 的候選集合共有 30 組，當預算達到或超過此數量時，改為遍歷全部組合。不同執行版本的預算差異見附錄~\ref{app:experimental-details}。

\clearpage
\begingroup
\footnotesize
\setstretch{1.1}
\setlength{\tabcolsep}{4pt}
\renewcommand{\arraystretch}{1.15}
\setlength{\LTleft}{0pt plus 1fill}
\setlength{\LTright}{0pt plus 1fill}
\begin{longtable}{@{}>{\raggedright\arraybackslash}p{\dimexpr 0.2100\linewidth-0.8400\tabcolsep\relax}>{\raggedright\arraybackslash}p{\dimexpr 0.3100\linewidth-1.2400\tabcolsep\relax}>{\raggedright\arraybackslash}p{\dimexpr 0.4800\linewidth-1.9200\tabcolsep\relax}@{}}
\caption{Baseline 傳統分類器之超參數候選集合}\label{tab:method-ml-search}\\
\toprule
\textbf{分類器} & \textbf{參數項目} & \textbf{超參數候選值} \\
\midrule
\endfirsthead
\multicolumn{3}{c}{表~\thetable（續）}\\
\toprule
\textbf{分類器} & \textbf{參數項目} & \textbf{超參數候選值} \\
\midrule
\endhead
\midrule
\multicolumn{3}{r}{續下頁}\\
\endfoot
\bottomrule
\endlastfoot
XGBoost & n\_\allowbreak{}estimators & 100、300、500、800、1000 \\
\addlinespace[1pt]
XGBoost & learning\_\allowbreak{}rate & 0.01、0.03、0.05、0.1、0.2 \\
\addlinespace[1pt]
XGBoost & max\_\allowbreak{}depth & 3、5、7、10 \\
\addlinespace[1pt]
XGBoost & min\_\allowbreak{}child\_\allowbreak{}weight & 1、3、5、20、50 \\
\addlinespace[1pt]
XGBoost & subsample & 0.7、0.85、1.0 \\
\addlinespace[1pt]
XGBoost & colsample\_\allowbreak{}bytree & 0.7、0.85、1.0 \\
\addlinespace[1pt]
XGBoost & gamma & 0.0、0.1、0.2 \\
\addlinespace[1pt]
XGBoost & reg\_\allowbreak{}alpha & 0.0、0.1、1.0 \\
\addlinespace[1pt]
XGBoost & reg\_\allowbreak{}lambda & 0.5、1.0、3.0、10.0 \\*
\addlinespace[1pt]
LightGBM & num\_\allowbreak{}leaves & 31、50、100、150 \\
\addlinespace[1pt]
LightGBM & max\_\allowbreak{}depth & 5、7、10、15、-1 \\
\addlinespace[1pt]
LightGBM & learning\_\allowbreak{}rate & 0.01、0.03、0.05、0.1、0.2 \\
\addlinespace[1pt]
LightGBM & n\_\allowbreak{}estimators & 100、200、500、800、1000 \\
\addlinespace[1pt]
LightGBM & min\_\allowbreak{}child\_\allowbreak{}samples & 20、50、100 \\
\addlinespace[1pt]
LightGBM & subsample & 0.7、0.85、1.0 \\
\addlinespace[1pt]
LightGBM & colsample\_\allowbreak{}bytree & 0.7、0.85、1.0 \\
\addlinespace[1pt]
LightGBM & reg\_\allowbreak{}alpha & 0.0、0.1、1.0 \\
\addlinespace[1pt]
LightGBM & reg\_\allowbreak{}lambda & 0.0、0.1、1.0、3.0 \\*
\addlinespace[1pt]
Random Forest & n\_\allowbreak{}estimators & 200、400、600、800 \\
\addlinespace[1pt]
Random Forest & max\_\allowbreak{}depth & None、8、12、16、20、24、32 \\
\addlinespace[1pt]
Random Forest & min\_\allowbreak{}samples\_\allowbreak{}leaf & 1、2、4、8、16 \\
\addlinespace[1pt]
Random Forest & min\_\allowbreak{}samples\_\allowbreak{}split & 2、4、8、16、32 \\
\addlinespace[1pt]
Random Forest & max\_\allowbreak{}features & sqrt、log2、0.2、0.35、0.5、0.65 \\*
\addlinespace[1pt]
Logistic Regression & C & 0.01、0.1、1.0、10.0、100.0 \\
\addlinespace[1pt]
Logistic Regression & solver & lbfgs、saga \\
\addlinespace[1pt]
Logistic Regression & max\_\allowbreak{}iter & 500、1000、2000 \\*
\end{longtable}
\endgroup

表中 None 表示 Random Forest 不限制樹深度，LightGBM 的深度值 -1 表示不設定深度上限；Random Forest 的小數型 \nolinkurl{max_features} 表示每次分裂考量的特徵比例。候選集合與實際執行版本之對應說明見附錄~\ref{app:experimental-details}。

\clearpage
\begingroup
\footnotesize
\setstretch{1.1}
\setlength{\tabcolsep}{4pt}
\renewcommand{\arraystretch}{1.15}
\setlength{\LTleft}{0pt plus 1fill}
\setlength{\LTright}{0pt plus 1fill}
\begin{longtable}{@{}>{\raggedright\arraybackslash}p{\dimexpr 0.2100\linewidth-0.8400\tabcolsep\relax}>{\raggedright\arraybackslash}p{\dimexpr 0.3100\linewidth-1.2400\tabcolsep\relax}>{\raggedright\arraybackslash}p{\dimexpr 0.4800\linewidth-1.9200\tabcolsep\relax}@{}}
\caption{Baseline TabM 之超參數候選集合}\label{tab:method-tabm-search}\\
\toprule
\textbf{分類器} & \textbf{參數項目} & \textbf{超參數候選值} \\
\midrule
\endfirsthead
\multicolumn{3}{c}{表~\thetable（續）}\\
\toprule
\textbf{分類器} & \textbf{參數項目} & \textbf{超參數候選值} \\
\midrule
\endhead
\midrule
\multicolumn{3}{r}{續下頁}\\
\endfoot
\bottomrule
\endlastfoot
TabM & k & 16、24、32 \\
\addlinespace[1pt]
TabM & n\_\allowbreak{}blocks & 2、3、4 \\
\addlinespace[1pt]
TabM & d\_\allowbreak{}block & 128、256、384、512 \\
\addlinespace[1pt]
TabM & dropout & 0.0、0.1、0.2、0.3 \\
\addlinespace[1pt]
TabM & lr & 0.0005、0.001、0.002、0.003 \\
\addlinespace[1pt]
TabM & weight\_\allowbreak{}decay & 1e-05、0.0001、0.0003、0.001 \\
\addlinespace[1pt]
TabM & use\_\allowbreak{}num\_\allowbreak{}embeddings & False、True \\
\addlinespace[1pt]
TabM & d\_\allowbreak{}embedding & 8、16 \\*
\end{longtable}
\endgroup

TabM 分別搜尋集成成員數、區塊數、區塊寬度、dropout、學習率、權重衰減與數值嵌入設定。未使用數值嵌入時，嵌入維度不參與該候選的有效設定。補充搜尋預算、訓練輪次上限、早停控制及 TabICL 固定參數，整理於附錄~\ref{app:experimental-details}。

\clearpage
\subsection{集成模型之固定設定}

實驗一的集成比較與實驗二沿用相同的基礎分類器與採樣設定，主要參數整理於表~\ref{tab:method-parameters}。模型未明示覆寫的參數使用固定套件版本之預設值。Baseline 的驗證選參與兩項集成比較的固定設定分別處理，其銜接與比較限制見附錄~\ref{app:experimental-details}。

\begingroup
\footnotesize
\setstretch{1.15}
\setlength{\tabcolsep}{4pt}
\renewcommand{\arraystretch}{1.25}
\setlength{\LTleft}{0pt plus 1fill}
\setlength{\LTright}{0pt plus 1fill}
\begin{longtable}{@{}>{\raggedright\arraybackslash}p{\dimexpr 0.2\linewidth-0.4000\tabcolsep\relax}>{\raggedright\arraybackslash}p{\dimexpr 0.8\linewidth-1.6000\tabcolsep\relax}@{}}
\caption{集成模型之主要方法與參數}\label{tab:method-parameters}\\
\toprule
\textbf{項目} & \textbf{固定設定} \\
\midrule
\endfirsthead
\multicolumn{2}{c}{表~\thetable（續）}\\
\toprule
\textbf{項目} & \textbf{固定設定} \\
\midrule
\endhead
\midrule
\multicolumn{2}{r}{續下頁}\\
\endfoot
\bottomrule
\endlastfoot
補值 & Old-fit 平均數補值；保留全空欄位的輸入維度 \\
\addlinespace[2pt]
縮放 & StandardScaler，僅擬合 Old-fit \\
\addlinespace[2pt]
基礎分類器 & XGBoost；{\normalfont objective=\allowbreak{}binary:logistic}、{\normalfont eval\_\allowbreak{}metric=\allowbreak{}auc}、{\normalfont tree\_\allowbreak{}method=\allowbreak{}hist}；不另加入類別權重；本部分未重新搜尋超參數，亦未啟用早停 \\
\addlinespace[2pt]
低採樣 & TomekLinks；{\normalfont sampling\_\allowbreak{}strategy=\allowbreak{}auto} \\
\addlinespace[2pt]
過採樣 & ADASYN；{\normalfont sampling\_\allowbreak{}strategy=\allowbreak{}auto}、{\normalfont n\_\allowbreak{}neighbors=\allowbreak{}5} \\
\addlinespace[2pt]
混合採樣 & SMOTEENN；SMOTE 的 {\normalfont k\_\allowbreak{}neighbors=\allowbreak{}5}；ENN 的 {\normalfont n\_\allowbreak{}neighbors=\allowbreak{}3}、{\normalfont kind\_\allowbreak{}sel=\allowbreak{}all} \\
\addlinespace[2pt]
MI & {\normalfont mutual\_\allowbreak{}info\_\allowbreak{}classif}；隨機種子隨實驗設定，依分數排序選取 \\
\addlinespace[2pt]
CART & DecisionTreeClassifier；{\normalfont class\_\allowbreak{}weight=\allowbreak{}balanced}，依特徵重要性排序選取 \\
\addlinespace[2pt]
RFE & 決策樹估計器；{\normalfont class\_\allowbreak{}weight=\allowbreak{}balanced}、{\normalfont step=\allowbreak{}1} \\
\addlinespace[2pt]
SHAP 排序 & LightGBM；{\normalfont n\_\allowbreak{}estimators=\allowbreak{}100}、{\normalfont class\_\allowbreak{}weight=\allowbreak{}balanced}；依 Old-fit 平均絕對貢獻值排序 \\
\addlinespace[2pt]
特徵保留 & 各方法分別保留 9、14 個；no-FS 使用 18 個 \\
\addlinespace[2pt]
靜態候選選擇 & 49 種混合子集，以 Calibration AP 排序；平手依固定方法名稱排序 \\
\addlinespace[2pt]
分類門檻 & Calibration F1-score；0.01–0.99、間隔 0.01；平手取最小門檻 \\
\addlinespace[2pt]
DCS／DES & OLA、LCA、KNORA-E、KNORA-U；{\normalfont k=\allowbreak{}7}、{\normalfont DFP=\allowbreak{}False}、{\normalfont with\_\allowbreak{}IH=\allowbreak{}False}；DCS 使用 {\normalfont best}，DES 使用 {\normalfont soft} \\*
\end{longtable}
\endgroup

2014 年分區種子固定為 42，所有模型種子與特徵條件均使用相同分區。模型、採樣與具有隨機性的特徵選取程序，使用 42、123、456、789、2024、7、21、84、168、999 共十個種子；可接收種子的元件依各次設定傳遞。Baseline 之超參數候選見第~\ref{sec:method-phase1-params}~節；補充訓練控制與版本說明見附錄~\ref{app:experimental-details}。


軟體環境、補充訓練控制與實驗紀錄規格詳見附錄~\ref{app:experimental-details}。

\section{實驗驗證準則}\label{sec:method-validation}

\subsection{訓練、驗證與測試程序}

各部分均先固定資料用途，再執行資料前處理與模型建立。基礎模型的學習、特徵選取及採樣僅使用允許的擬合資料；DSEL 用於動態局部能力估計；Calibration 用於靜態候選及分類門檻決定。測試標籤不參與上述決策。

使用同一 Calibration 比較候選並決定門檻，屬於本研究的模型選擇程序，因此 Calibration 上的最佳分數不作為最終泛化效能。正式報告以決策固定後的 Test 表現為準，並保留各方法與各特徵條件的結果，不只呈現測試期表現最高者。

\subsection{分類效能評估指標}

本研究將 failed 定義為正類。TP 為實際破產且預測為破產的紀錄數；FN 為實際破產但預測為存續的紀錄數；FP 為實際存續但預測為破產的紀錄數；TN 為實際存續且預測為存續的紀錄數。

\textbf{一、AUC}

ROC 曲線以真陽性率（True Positive Rate, TPR）為縱軸，假陽性率（False Positive Rate, FPR）為橫軸，定義為：

\begin{equation}\label{eq:method-roc}
\mathrm{TPR}=\frac{TP}{TP+FN},\qquad
\mathrm{FPR}=\frac{FP}{FP+TN}.
\end{equation}

AUC 為 ROC 曲線下的面積，用於衡量正負類樣本的排序區辨能力，數值越高表示整體區辨表現越佳。AUC 由預測機率計算，不依賴本研究所選的單一分類門檻；惟其無法單獨反映實際門檻下的破產樣本辨識情形，故另以 F1-score 與 Recall 補充。

\textbf{二、F1-score}

Precision 表示預測為破產的紀錄中，實際破產者所占比例；Recall 表示實際破產紀錄中被正確辨識的比例。兩者與 F1-score 定義如下：

\begin{equation}\label{eq:recall}
\mathrm{Precision}=\frac{TP}{TP+FP},\qquad
\mathrm{Recall}=\frac{TP}{TP+FN},
\end{equation}

\begin{equation}\label{eq:method-f1}
F1=\frac{2\,\mathrm{Precision}\,\mathrm{Recall}}
{\mathrm{Precision}+\mathrm{Recall}}
=\frac{2TP}{2TP+FP+FN}.
\end{equation}

F1-score 同時反映正類預測的精確程度與涵蓋程度。本研究以各方法在 Calibration 決定的門檻計算測試 F1-score；計算遇到未定義的零分母情形時，依統一的評估規則記為零。

\textbf{三、Recall}

Recall 與 TPR 相同，表示實際破產紀錄中被正確辨識的比例。較低的 Recall 代表較多破產紀錄被漏判；惟較高 Recall 亦可能伴隨較多誤報，因此須與 F1-score 及 AUC 共同解讀，不僅依 Recall 判定方法優劣。

\textbf{四、Average Precision}

AP 用於實驗一之集成比較與實驗二的靜態候選選擇，依不同門檻的 Precision 與 Recall 計算：

\begin{equation}\label{eq:method-ap}
\mathrm{AP}=\sum_n(R_n-R_{n-1})P_n,
\end{equation}

其中，\(P_n\) 與 \(R_n\) 分別為第 \(n\) 個操作點的 Precision 與 Recall。本研究以 Recall 增量加權 Precision；此計算不同於對 PR 曲線直接進行梯形積分，定義見 scikit-learn AP 文件\autocite{sklearn-ap-docs}。AP 用於驗證選模，最終主文仍以 AUC、F1-score 與 Recall 為主要報告指標。

\subsection{結果彙整與模型比較}

實驗一之集成比較與實驗二分別計算 2015 年至 2018 年的 pooled 與逐年指標。Pooled 指將四個測試年度的標籤與預測合併後計算一次指標，與先計算四個年度指標再取平均不同；本研究不混用兩種口徑。逐年結果用於觀察不同年度的分類表現，所有年度沿用驗證程序固定的模型與門檻。

跨種子結果以平均值與種子標準差呈現。十個種子使用相同資料來源、驗證分區與測試樣本，標準差主要描述本次設定下的程序隨機變異，不能視為十份獨立資料的抽樣不確定性，亦不直接作為統計顯著性證據。公司層級的配對檢定另依第~\ref{sec:method-statistics}~節執行。

實驗一分別呈現 Baseline 的分類器與訓練期間比較，以及固定年份的模型池與集成方法比較；實驗二依相同種子、方法與測試樣本，對照未進行 FS 的集成結果。固定方法的配對比較與重新選擇靜態成員的政策比較分別呈現。研究結論依此比較範圍解釋，不將跨協定分數差直接歸因於集成或特徵選取，亦不由單一資料集的結果推論跨市場或未來事件預測能力。

\subsection{公司群聚配對統計檢定}\label{sec:method-statistics}

為評估方法間差異的不確定性，本研究針對固定年份集成結果，補作以公司為重抽單位的配對群聚 bootstrap。此分析在既有實驗完成後制定，屬事後探索性分析；所有分類器、特徵設定、靜態選模規則與分類門檻均使用已保存的驗證決策，不因檢定結果調整模型。Bootstrap 的一般原理見 Efron 與 Tibshirani\autocite{efron1986bootstrap}；群聚 ROC 資料須考量同一群聚內相關性的原因，見 Obuchowski\autocite{obuchowski1997clustered}。本研究依公司年度資料結構設計下列配對重抽程序。

\subsubsection{比較對象與估計量}

實驗一以完整六模型等權集成 Equal6 為共同基準，分別比較 Old3、New3、驗證選定靜態集成（以下簡稱 V-static）、DCS-OLA、DCS-LCA、KNORA-E 及 KNORA-U，共七組比較。V-static 在各種子內依 Calibration AP 選擇混合成員，不使用 Test 排序結果選擇靜態方案。Baseline 的 15 組年份邊界共享訓練與測試資料，且與固定年份集成的驗證制度不同，因此不將它們當成獨立重複進行顯著性檢定，也不與集成結果直接作跨協定檢定。

實驗二將 MI、CART、SHAP 與 RFE 的兩種保留比例，各自與相同方法的 no-FS 結果配對。比較方法為 Equal6、V-static、DCS-OLA、DCS-LCA、KNORA-E 與 KNORA-U，因此共有 \(8\times6=48\) 組比較。V-static 允許 FS 後依原有 Calibration 規則重新選擇成員，故該比較評估完整選模程序的變化；其他方法保持方法名稱與規則一致，但仍重新訓練模型並以 Calibration 固定各自門檻。

每組比較均計算 AUC、F1-score 及 Recall，合計 \((7+48)\times3=165\) 項檢定。若 \(M\) 為任一指標，\(S=10\) 為種子數，方法 A 相對方法 B 的差值定義為：
\begin{equation}\label{eq:stats-difference}
\widehat{\Delta}_M=\frac{1}{S}\sum_{s=1}^{S}
\left[M(y,p_{A,s};\tau_{A,s})-M(y,p_{B,s};\tau_{B,s})\right].
\end{equation}
其中，\(p\) 為模型預測機率，\(\tau\) 為驗證程序固定的門檻；AUC 不使用門檻。此估計量先計算各種子的 pooled 指標，再平均其差值，不將十個種子的機率平均後另建一個預測模型。

\subsubsection{重抽、信賴區間與多重比較}

每次 bootstrap 從 Test 所含的 \(C\) 家公司中等機率、有放回抽取 \(C\) 次。公司被抽中後，保留該公司在 2015--2018 年的全部測試紀錄，並將同一組公司抽取次數同時套用於所有種子、方法及 FS 設定。如此可保留同公司跨年的相依性及方法間的配對關係；各次重抽均重新計算指標，但不重新訓練或調整門檻。

本研究使用 49,999 次重抽，重抽種子固定為 20261027。由差值重抽分布的第 2.5 與第 97.5 百分位數建立逐項 95\% 信賴區間。以 \(\Delta_M^{*(b)}\) 表示第 \(b\) 次重抽差值，針對虛無假設 \(H_0:\Delta_M=0\)，以置中重抽誤差分布計算雙尾近似 \(p\) 值：
\begin{equation}\label{eq:stats-bootstrap-p}
p_M=\frac{1+\sum_{b=1}^{B}\mathbf{1}\!\left\{
\left|\Delta_M^{*(b)}-\widehat{\Delta}_M\right|\geq
\left|\widehat{\Delta}_M\right|\right\}}{B+1},
\qquad B=49{,}999.
\end{equation}
此為有限次 Monte Carlo 計算的 bootstrap 近似檢定，並非精確置換檢定；加一修正使報告值不為零，最小可報值為 0.00002。

為避免只呈現有利的比較，兩項實驗、所有方法與三項指標的 165 個 \(p\) 值共同採 Holm 法\autocite{holm1979multiple} 校正，整體顯著水準設為 0.05，校正後數值記為 \(p_H\)。區間為未作多重比較校正的逐項區間，因此即使區間不包含零，仍可能未通過 Holm 校正；顯著性判斷以 \(p_H\) 為準。未拒絕虛無假設不表示兩方法等效。

此推論以既定的訓練模型、資料切分、驗證決策與十個種子為條件，描述測試公司組成改變時的差異，不涵蓋重新抽取訓練公司的不確定性。方法亦假設不同公司可作為重抽單位，未另外建模公司間產業關聯或共同年度衝擊，故不能據此推論其他景氣期間、未見公司或未來事件預測的普遍優越性。檢定結果於第四章呈現，完整比較列於附錄。
\endgroup
